
\documentclass[12pt,a4paper]{report}

% -------------------- PACKAGES --------------------
\usepackage[utf8]{inputenc}
\usepackage[T1]{fontenc}
\usepackage{geometry}
\usepackage{graphicx}
\usepackage{float}
\usepackage{booktabs}
\usepackage{array}
\usepackage{longtable}
\usepackage{hyperref}
\usepackage{xcolor}
\usepackage{caption}
\usepackage{fancyhdr}
\usepackage{enumitem}
\usepackage{amsmath}
\usepackage{listings}

\geometry{margin=1in}
\setlength{\parskip}{6pt}
\setlength{\parindent}{0pt}

\hypersetup{
    colorlinks=true,
    linkcolor=blue,
    urlcolor=blue,
    citecolor=blue
}

\pagestyle{fancy}
\fancyhf{}
\lhead{Activity 2 -- Research Report}
\rhead{Department of MCA}
\cfoot{\thepage}

% -------------------- DOCUMENT --------------------
\begin{document}

% -------------------- TITLE PAGE --------------------
\begin{titlepage}
\centering

{\Large \textbf{[NAME OF THE INSTITUTION]}}\\[0.3cm]
{\large Department}\\[1.5cm]

{\Huge \textbf{ACTIVITY 2}}\\[0.4cm]
{\Large Student Research Report}\\[1.2cm]

\hrule
\vspace{0.5cm}

{\LARGE \textbf{[TITLE OF THE PROJECT]}}\\[0.4cm]
{\large Team Name: [TEAM NAME]}

\vspace{0.5cm}
\hrule
\vspace{1cm}

\textbf{Submitted by}\\[0.3cm]

\begin{tabular}{ll}
Student 1: & [Name and USN] \\
Student 2: & [Name and USN] \\
Student 3: & [Name and USN] \\
Student 4: & [Name and USN] \\
\end{tabular}

\vfill

\textbf{Course:} [Course Name]\\
\textbf{Faculty:} [Faculty Name]\\
\textbf{Academic Year:} [Academic Year]\\[0.5cm]

\textbf{Submission Date:} [DD/MM/YYYY]

\end{titlepage}

\pagenumbering{roman}
\tableofcontents
\clearpage
\pagenumbering{arabic}

% ==================================================
\chapter{Introduction}

\textbf{Instructions:}
\begin{itemize}
    \item Introduce the research topic and its background.
    \item Explain why the problem is important in the real world.
    \item Describe the existing challenges and the motivation
    for selecting this project.
    \item Clearly define the problem statement.
    \item State the objectives and scope of the project.
\end{itemize}

\section{Background}

[Introduce the topic, its applications, and its importance.]

\section{Problem Statement}

[Clearly explain the specific problem that the project
aims to solve. Identify the limitations of existing
approaches and the problem your team intends to address.]

\section{Objectives}

\begin{enumerate}
    \item [Objective 1]
    \item [Objective 2]
    \item [Objective 3]
\end{enumerate}

\section{Scope of the Project}

[Describe what the project covers, its intended users,
and its major limitations or boundaries.]

% ==================================================
\chapter{Literature Survey}

\textbf{Instructions:}
\begin{itemize}
    \item Review at least 8 relevant research papers
    (or the number specified by the faculty).
    \item Prefer peer-reviewed journals, conference papers,
    and reliable academic publications.
    \item For every paper, mention its title, publication
    venue, publication year, methodology, and major results.
    \item Explain the limitations or research gaps identified.
    \item Cite every paper in the text and include its full
    reference in the References section.
    \item Do not copy abstracts directly. Summarize the
    research in your own words.
\end{itemize}

\section{Comparison of Existing Research}

\small
\begin{longtable}{
|p{3.2cm}|p{2.0cm}|p{1.0cm}|p{3.0cm}|p{3.0cm}|
}
\caption{Comparative Analysis of Research Papers}\\
\hline
\textbf{Paper Title} &
\textbf{Publication} &
\textbf{Year} &
\textbf{Methodology} &
\textbf{Results / Findings} \\
\hline
\endfirsthead

\hline
\textbf{Paper Title} &
\textbf{Publication} &
\textbf{Year} &
\textbf{Methodology} &
\textbf{Results / Findings} \\
\hline
\endhead

[Paper 1] & [Journal/Conference] & [Year] &
[Algorithm or approach used] &
[Main findings and reported performance] \\
\hline

[Paper 2] & [Journal/Conference] & [Year] &
[Algorithm or approach used] &
[Main findings and reported performance] \\
\hline

[Paper 3] & [Journal/Conference] & [Year] &
[Algorithm or approach used] &
[Main findings and reported performance] \\
\hline

[Paper 4] & [Journal/Conference] & [Year] &
[Algorithm or approach used] &
[Main findings and reported performance] \\
\hline

[Paper 5] & [Journal/Conference] & [Year] &
[Algorithm or approach used] &
[Main findings and reported performance] \\
\hline

\end{longtable}
\normalsize

\section{Research Gap}

[Identify what is missing, insufficient, or unresolved
in the existing studies. Explain how your project attempts
to address at least one of these gaps.]

% ==================================================
\chapter{Methodology}

\textbf{Instructions:}
\begin{itemize}
    \item Explain the complete workflow of your proposed solution.
    \item Describe the dataset, its source, and its characteristics.
    \item Explain data preprocessing and feature selection, if applicable.
    \item Describe the algorithms, models, tools, and technologies used.
    \item Include a system architecture or workflow diagram.
    \item Explain the training, validation, and testing procedure,
    wherever applicable.
    \item The methodology must be detailed enough for another
    student to understand and reproduce your work.
\end{itemize}

\section{Proposed Approach}

[Describe the proposed solution and how it addresses
the problem statement.]

\section{Dataset Description}

\begin{itemize}
    \item Dataset name and source: [Provide details]
    \item Number of samples: [Enter count]
    \item Number of features: [Enter count]
    \item Target classes: [List classes, if applicable]
    \item Preprocessing steps: [Describe steps]
\end{itemize}

\section{Algorithms and Technologies}

[Describe the algorithms, libraries, programming languages,
frameworks, and tools used in the implementation.]

\section{System Workflow}

% Upload your workflow diagram to the Overleaf project
% and replace workflow.png with its filename.

\begin{figure}[H]
\centering
\fbox{\parbox{0.8\textwidth}{\centering
Insert your system architecture or workflow diagram here.}}
% \includegraphics[width=0.85\textwidth]{workflow.png}
\caption{Proposed System Workflow}
\label{fig:workflow}
\end{figure}

% ==================================================
\chapter{Experimental Details and Results}

\textbf{Instructions:}
\begin{itemize}
    \item Mention the hardware and software environment.
    \item Describe the experimental setup and parameter settings.
    \item Specify the training and testing split, where applicable.
    \item Report the actual results obtained from your implementation.
    \item Use suitable evaluation metrics for your problem.
    \item Include screenshots of actual outputs and plots.
    \item For classification projects, include the confusion matrix
    and explain the values.
    \item Compare your results with baseline models or published
    results only when the comparison is valid.
\end{itemize}

\section{Experimental Environment}

\begin{itemize}
    \item Programming language: [Python / Java / Other]
    \item Development environment: [IDE / Notebook / Platform]
    \item Libraries and frameworks: [List]
    \item Hardware configuration: [CPU / GPU / RAM]
    \item Dataset split: [Training / Validation / Testing]
\end{itemize}

\section{Evaluation Metrics}

[Define the evaluation metrics appropriate to your project.
For a classification task, these may include accuracy,
precision, recall, F1-score, and ROC-AUC.]

\section{Experimental Results}

\begin{table}[H]
\centering
\caption{Performance of the Proposed Model}
\begin{tabular}{|l|c|}
\hline
\textbf{Evaluation Metric} & \textbf{Obtained Value} \\
\hline
Accuracy & [Value or N/A] \\
\hline
Precision & [Value or N/A] \\
\hline
Recall & [Value or N/A] \\
\hline
F1-score & [Value or N/A] \\
\hline
Other relevant metric & [Value or N/A] \\
\hline
\end{tabular}
\end{table}

[Interpret the results. Explain what the metrics indicate
about the performance of your approach. Use N/A for metrics
that are not applicable and justify the choice of metrics.]

\section{Confusion Matrix}

\textbf{Instructions:}
Include the confusion matrix generated by your actual model.
Label the axes and classes clearly. Explain true positives,
true negatives, false positives, and false negatives
where applicable. For multiclass classification, explain
the major misclassifications between classes.

\begin{figure}[H]
\centering
\fbox{\parbox{0.8\textwidth}{\centering
Insert your actual confusion matrix screenshot here.}}
% \includegraphics[width=0.65\textwidth]{confusion_matrix.png}
\caption{Confusion Matrix of the Proposed Model}
\label{fig:confusion}
\end{figure}

\section{Output Screenshots}

\textbf{Instructions:}
Include screenshots showing actual execution of your project.
Screenshots should be readable and relevant. Add a caption
and briefly explain each output.

\begin{figure}[H]
\centering
\fbox{\parbox{0.8\textwidth}{\centering
Insert your actual project output screenshot here.}}
% \includegraphics[width=0.85\textwidth]{results.png}
\caption{Sample Output of the Proposed System}
\label{fig:output}
\end{figure}

% ==================================================
\chapter{Error Analysis}

\textbf{Instructions:}
\begin{itemize}
    \item Identify incorrect predictions, failures, or unexpected outputs.
    \item Explain why these errors may have occurred.
    \item Discuss the impact of data quality, class imbalance,
    model limitations, or parameter selection, as relevant.
    \item Provide examples or screenshots of failure cases, if available.
    \item Suggest practical ways to reduce the errors.
    \item If no errors were observed in a particular test,
    state the tests performed and their limitations.
\end{itemize}

\section{Observed Errors}

[Describe the errors or failure cases observed during
testing and provide evidence where possible.]

\section{Possible Causes}

[Discuss the possible reasons for the errors, supported
by experimental observations.]

\section{Improvement Strategies}

[Explain how the model or system could be improved.
Distinguish implemented improvements from future suggestions.]

% ==================================================
\chapter{Project Proof and Implementation Evidence}

\textbf{Instructions:}
Submit verifiable evidence that the project has been
implemented. Include the repository link and screenshots
of the running application, if deployed.

\section{GitHub Repository}

\textbf{Repository URL:}
\url{https://github.com/your-username/your-repository}

Replace the sample URL with your team's actual repository.
The repository should contain the source code, README,
setup instructions, and other essential project files.

\section{Deployment Evidence}

\textbf{Deployment URL (if applicable):}
\url{https://your-deployment-url.example}

If the project has not been deployed, state
``Not deployed'' and explain how it can be executed locally.

\section{Deployed Application Screenshot}

\begin{figure}[H]
\centering
\fbox{\parbox{0.8\textwidth}{\centering
Insert a screenshot of your deployed application here.}}
% \includegraphics[width=0.85\textwidth]{deployment.png}
\caption{Deployed Project Interface}
\label{fig:deployment}
\end{figure}

[Explain what the screenshot demonstrates. Do not include
passwords, API keys, private tokens, or other credentials.]

% ==================================================
\chapter{Conclusion and Future Work}

\textbf{Instructions:}
\begin{itemize}
    \item Summarize the problem addressed and the proposed solution.
    \item State the main findings based on actual results.
    \item Explain whether the project objectives were achieved.
    \item Mention important limitations.
    \item Suggest realistic directions for future improvement.
\end{itemize}

\section{Conclusion}

[Summarize the project, its implementation, and the
most important findings. Support claims using your results.]

\section{Future Work}

[Describe possible extensions, additional experiments,
or improvements that could be explored in future.]

% ==================================================
\chapter{Acknowledgement of AI Usage}

\textbf{Instructions:}
\begin{itemize}
    \item Disclose whether AI tools were used during the project.
    \item Mention the tool names and the purpose of their use.
    \item Explain which outputs were reviewed, tested, or modified
    by the team.
    \item Clearly identify any AI-generated code, text, or ideas
    that materially contributed to the work, as required by
    the course or institutional policy.
    \item Students remain responsible for the correctness of
    the code, results, citations, and final report.
\end{itemize}

\section{AI Usage Statement}

[Complete one of the following statements and customize it.]

\textbf{If AI tools were used:}

We used [AI tool names] for [specific purposes, such as
brainstorming, code assistance, debugging, or language editing].
The team reviewed and validated the relevant outputs.
All experiments and reported results were checked against
our actual implementation. The team takes responsibility
for the final submission.

\textbf{If AI tools were not used:}

No generative AI tools were used in preparing this report
or developing the project.

[Modify the statement to reflect actual usage. Do not
claim that AI was or was not used inaccurately.]

% ==================================================
\begin{thebibliography}{99}

\bibitem{ref1}
[Author(s)], ``[Title of research paper],''
\textit{[Journal or Conference Name]}, [Volume, Issue],
pp. [Page numbers], [Year].
DOI: [DOI or URL].

\bibitem{ref2}
[Author(s)], ``[Title of research paper],''
\textit{[Journal or Conference Name]}, [Year].
DOI: [DOI or URL].

\bibitem{ref3}
[Author(s)], ``[Title of research paper],''
\textit{[Journal or Conference Name]}, [Year].
DOI: [DOI or URL].

\bibitem{ref4}
[Author(s)], ``[Title of research paper],''
\textit{[Journal or Conference Name]}, [Year].
DOI: [DOI or URL].

\bibitem{ref5}
[Author(s)], ``[Title of research paper],''
\textit{[Journal or Conference Name]}, [Year].
DOI: [DOI or URL].

\end{thebibliography}

\end{document}